\section{Selecting candidates from a scan}
\label{sec:selection}

Sections~\ref{sec:budget} and~\ref{sec:combinatorial} priced the scan. This section asks what to
do with its output. A broad scan produces a large number of local significances, one per bin per
spectrum, and something has to decide which of them is worth following up. Three rules are
compared here on identical footing: take the largest (\emph{argmax}), take everything above a
fixed threshold, or let the data set the cut through a false-discovery-rate criterion.

The setting is deliberately minimal: $n = 30\,000$ bins, of which $n-1$ are background with $z_i
\sim \mathcal{N}(0,1)$, and at most one contains a signal with $z \sim \mathcal{N}(\mu, 1)$.
Whatever passes stage 1 is unblinded in an independent sample B and confirmed if $z_B > 3$. This
is the two-stage structure of Section~\ref{sec:ab} stripped to its statistical core.

\subsection{The background maximum is not Gaussian}

The first thing to get right is the null distribution of $\max_i z_i$. Two limits are routinely
confused. Running more pseudo-experiments makes the histogram converge to the \emph{exact} density
\begin{equation}
  f_M(x) \;=\; n\,\phi(x)\,\Phi(x)^{n-1},
  \qquad
  F_M(x) = \Phi(x)^n,
  \qquad
  x_q = \Phi^{-1}\!\left(q^{1/n}\right),
  \label{eq:maxexact}
\end{equation}
which is exact for every $n$. Increasing $n$ makes the \emph{rescaled} $f_M$ converge to a Gumbel
law~\cite{fishertippett,gumbel}, and that convergence is only logarithmic, with relative error
$\mathcal{O}(1/2\ln n)$.

At $n = 30\,000$ the difference is measurable, and $10^4$ pseudo-experiments are enough to resolve
it (Table~\ref{tab:maxvalidation}). The exact form is confirmed; the asymptotic Gumbel is rejected
at $p = 2\times10^{-15}$ and a best-fit Gaussian at $p = 6\times10^{-30}$.

\begin{table}[htbp]
\centering
\caption{Distribution of the maximum of $n=30\,000$ standard normals, from $10^4$
pseudo-experiments, against the exact law and its Gumbel limit
($a_n = 4.005$, $b_n = 0.220$).}
\label{tab:maxvalidation}
\begin{tabular}{lrr}
\toprule
 & simulated & Gumbel asymptotic \\
\midrule
mean & 4.110 & 4.132 \\
standard deviation & 0.284 & 0.282 \\
median & 4.072 & 4.086 \\
skewness & 0.89 & 1.14 \\
\midrule
KS vs exact $\Phi(x)^n$ & \multicolumn{2}{r}{$D = 0.0077$, $p = 0.59$ \quad\checkmark} \\
KS vs Gumbel$(a_n,b_n)$ & \multicolumn{2}{r}{$D = 0.041$, $p = 2\times10^{-15}$} \\
KS vs best-fit Gaussian & \multicolumn{2}{r}{$D = 0.058$, $p = 6\times10^{-30}$} \\
\bottomrule
\end{tabular}
\end{table}

\begin{figure}[htbp]
\centering
\includegraphics[width=0.95\textwidth]{max_of_gaussians_light.png}
\caption{Left: one pseudo-experiment of $30\,000$ Gaussian draws, with its maximum marked. Right: the
distribution of $10^4$ such maxima, with the exact density of Eq.~\eqref{eq:maxexact} and the
asymptotic Gumbel overlaid. The distribution is strongly right-skewed, so a symmetric error bar on
a maximum is never appropriate.}
\label{fig:maxgauss}
\end{figure}

The practical rules that follow are short. Gumbel is fine for intuition. For any actual number,
whether a threshold, a $p$-value or a quantile, use $F_M = \Phi^n$, which costs nothing and is
right. Never model a maximum as Gaussian. This also connects back to Section~\ref{sec:budget}: the
familiar trials factor is the leading term of $1 - \Phi(z)^{N_{\mathrm{eff}}}$, and in the far
tail, where global $p$-values live, a few-percent Gumbel error is not something one wants between
oneself and a discovery claim.

\subsection{A real signal often loses the scan}

The probability that the signal bin \emph{is} the maximum, $p(\mu) = \int
\phi(s-\mu)\,\Phi(s)^{n-1}\,\mathrm{d}s$, is sobering. A genuine $4\sigma$ signal wins the scan
only 46\,\% of the time, because it is racing a distribution centred at $4.11$ with a spread of
$0.28$. The two-stage confirmation probability follows as
$p_{\mathrm{conf}} = p(\mu)\,\Phi(\mu - 3) + [1-p(\mu)]\,\Phi(-3)$ (Table~\ref{tab:rules}).

The second stage is nearly free. For $\mu \ge 5\sigma$ the loss between ``wins the scan'' and
``confirmed'' is under two percentage points, because a signal strong enough to win stage 1 almost
always exceeds $3\sigma$ in stage 2. \textbf{Stage 2 also carries no trials factor}, since the bin
was chosen by stage 1. The false-alarm rate per opening is $1-\Phi(3) = 1.35\times10^{-3}$, or
$3.0\sigma$ \emph{global}, whatever the size of the scan that selected it.

\subsection{A threshold dominates the argmax}

Taking the maximum throws information away: it ranks bins and then discards everything about their
absolute size. The alternative is to select every bin above a fixed threshold $t$. Under
background-only the number selected is Poisson with mean $\lambda(t) = (n-1)\,[1-\Phi(t)]$, so the
natural operating point is $t^\star = \Phi^{-1}(1 - 1/(n-1)) = 3.99$, where $\lambda = 1$. That is
exactly the argmax's fake budget of one candidate per experiment, and at that matched budget the
threshold rule wins at every signal strength (Table~\ref{tab:rules}).

\begin{table}[htbp]
\centering
\caption{Probability that a signal of strength $\mu$ is confirmed in sample B, for the three
stage-1 rules, all at the same background-only budget of one candidate per experiment
($3.0\sigma$ global false-alarm rate). The threshold operating point is $t^\star = 3.99$ and the
matched false-discovery rate is $q^\star = 0.381$.}
\label{tab:rules}
\begin{tabular}{lrrrr}
\toprule
$\mu$ & wins the scan & argmax confirmed & BH at $q^\star$ & threshold at $t^\star$ \\
\midrule
$3\sigma$ & 14.2\,\% & 7.2\,\% & 7.5\,\% & \textbf{8.1\,\%} \\
$4\sigma$ & 45.8\,\% & 38.6\,\% & 40.3\,\% & \textbf{42.5\,\%} \\
$5\sigma$ & 80.3\,\% & 78.5\,\% & 80.6\,\% & \textbf{82.5\,\%} \\
$6\sigma$ & 96.4\,\% & 96.3\,\% & 97.2\,\% & \textbf{97.7\,\%} \\
\bottomrule
\end{tabular}
\end{table}

Putting both rules on a common ROC, with the average number of false confirmations on one axis and
the signal confirmation probability on the other, makes the structure clear
(Figure~\ref{fig:roc}). The threshold rule traces a curve as $t$ varies. The argmax is a
\emph{single point}, because it always selects exactly one bin, and that point lies strictly
inside the curve at every $\mu$. It is dominated, not merely different.

One caveat took some care to get right. The bare ratio of power to false positives has no interior
maximum; it increases monotonically with $t$ and simply says ``cut harder''. The well-posed
question is where the ROC tangent has a given slope, which amounts to saying what a false positive
costs relative to a discovery. Once that cost is specified, the optimum has a closed form.

\begin{figure}[htbp]
\centering
\includegraphics[width=0.95\textwidth]{roc_threshold_vs_argmax.png}
\caption{Left: the ROC of the threshold rule (one curve per signal strength) with the argmax rule as a
single point (diamonds), at $\mu = 3, 4, 5, 6\sigma$. The argmax point lies inside the curve at
every $\mu$, so at the same false-positive budget a threshold confirms more signal. Right: the
bare power-per-false-positive ratio, which never turns over. An objective has to specify what a
false positive costs.}
\label{fig:roc}
\end{figure}

\subsection{Letting the data set the cut: Benjamini--Hochberg}

The third rule replaces a fixed bar with a false-discovery-rate
criterion~\cite{benjaminihochberg}. Convert each bin to a one-sided $p$-value, sort, and reject
the $K(q)$ smallest, where $K(q) = \max\{k : p_{(k)} \le kq/n\}$. The dial is now $q$, the nominal
FDR. Since $K$ is a random variable depending on the whole $p$-vector, this has to be evaluated by
Monte Carlo.

Two properties matter here. First, under the global null the probability that BH makes at least one rejection is exactly $q$.
This follows from the identity $P(\sup_t F_n(t)/t \ge 1/q) = q$ for the uniform empirical
distribution, usually attributed to Daniels~\cite{daniels}, and our Monte Carlo reproduces it ($q
= 0.05 \to 0.0505$, $0.2 \to 0.1990$, $0.5 \to 0.5008$). The \emph{number} of rejections is not $q$: BH is
bursty, and past $q^\star$ multi-rejection events dominate. Second, the realised cut is not fixed.
In a pseudo-experiment where BH rejects $K$ bins the weakest accepted bin sits at
$z_{\mathrm{cut}} = \Phi^{-1}(1-p_{(K)})$, which fluctuates from experiment to experiment.

Matching BH to the same fake budget as the other two rules, $E[K \mid H_0] = 1$, gives $q^\star =
0.381$, and it then agrees on both false-alarm measures ($1.34\times10^{-3}$ expected false
confirmations and $P(\ge1) = 1.336\times10^{-3}$, which is $3.0\sigma$ global). Its effective
single-bin bar is $z_1(q^\star) = 4.21$, \emph{stricter} than $t^\star = 3.99$.

BH lands strictly between the other two rules (Table~\ref{tab:rules}, Figure~\ref{fig:rocbh}). It
beats the argmax for the same reason the threshold does: it never discards the signal candidate
merely because some background bin happened to be larger, and the step-up lets a signal ride along
at rank 2 or below. It loses to the fixed threshold because keeping $E[K\mid H_0] = 1$ while
spending part of that budget on rank-$2+$ clumps forces the single-bin bar up to $4.21$, and those
extra selections are pure waste. They are background bins, and the signal, normally rank 1, gains
nothing from them. The effect is specific to large $q$; at small $q$ BH is essentially the rank-1
bar and the three rules converge.

\begin{figure}[htbp]
\centering
\includegraphics[width=0.95\textwidth]{roc_bh_vs_threshold.png}
\caption{Left: all three stage-1 rules on the ROC of Figure~\ref{fig:roc}. BH runs just inside the
fixed-threshold curve and outside the argmax point. Right: the realised false discovery rate
before and after the stage-B confirmation. The confirmation step drives it to near zero regardless
of the nominal $q$.}
\label{fig:rocbh}
\end{figure}

\subsection{The same three rules with an imperfect estimator}
\label{sec:selnoise}

Everything above assumes the estimator returns $z \sim \mathcal{N}(0,1)$ for every background bin.
Section~\ref{sec:abnoise} argues that a network-driven scan does not have that estimator, and the
comparison is worth redoing with the same defect model, because the three rules are not equally
exposed to it. The exposure is structural rather than numerical: argmax and BH are
\textbf{rank-based}, so an artefact anywhere in the scan changes the decision about every other
bin, while a fixed threshold is \textbf{absolute} and an artefact at 1\,TeV says nothing about the
bin at 3\,TeV.

Table~\ref{tab:selnoise} redoes Table~\ref{tab:rules} at $\mu = 5\sigma$, with every rule tuned to
the \emph{same} $3.0\sigma$-global false-alarm budget under each contamination. The
$\epsilon = 0$ row reproduces the perfect-estimator result, which fixes the comparison.

\begin{table}[htbp]
\centering
\caption{Probability that a $5\sigma$ signal is confirmed, for the three stage-1 rules, all held
at $1.35\times10^{-3}$ expected false confirmations. \textsc{glitch} and \textsc{bias} are the
incoherent and coherent defect classes of Section~\ref{sec:abnoise}. ``n/a'' means the rule cannot
be tuned to that budget at all. Monte-Carlo error on the BH column is $\pm0.35$ points; the other
two are analytic.}
\label{tab:selnoise}
\begin{tabular}{lrrrrr}
\toprule
estimator & argmax & threshold & BH & thr $-$ argmax & thr $-$ BH \\
\midrule
perfect (Table~\ref{tab:rules}) & 78.5\,\% & \textbf{82.5\,\%} & 80.9\,\% & $+4.0$ & $+1.6$ \\
\textsc{glitch} $\epsilon = 10^{-4}$ & 73.2\,\% & \textbf{81.2\,\%} & 80.1\,\% & $+8.0$ & $+1.1$ \\
\textsc{glitch} $\epsilon = 10^{-3}$ & 43.8\,\% & \textbf{57.3\,\%} & 57.2\,\% & $+13.6$ & $+0.1$ \\
\textsc{glitch} $\epsilon = 10^{-2}$ & 6.2\,\% & 5.8\,\% & 6.0\,\% & $-0.4$ & $-0.2$ \\
\textsc{bias}, any $\epsilon$ & \emph{untunable} & $\sim 0\,\%$ & \emph{unreachable} & --- & --- \\
\bottomrule
\end{tabular}
\end{table}

\paragraph{The threshold's margin over the argmax widens, from $+4$ to $+13.6$ points.} The cause
is visible in a single diagnostic: the probability that the argmax \emph{is} the signal falls from
$80.3\,\%$ to $74.9\,\%$ to $44.8\,\%$ as $\epsilon$ runs from $0$ to $10^{-3}$. At $\epsilon =
10^{-3}$ some thirty artefact bins are scattered through a $30\,000$-bin scan, and any one of them
outranking the signal removes it completely. A rank-1 rule has no graceful degradation, whereas a
threshold pays only through the extra candidates it forwards.

\paragraph{Contamination strips BH of the adaptivity that made it attractive.} BH converges onto
the fixed threshold ($+1.6 \to +1.1 \to +0.1$), and the reason is mechanical: to hold the
false-alarm budget it has to be run at ever smaller $q$, and BH at small $q$ simply \emph{is} a
threshold at $\Phi^{-1}(1-q/n)$. One keeps the threshold's performance while losing the ability to
interpret the dial.

\paragraph{And that dial becomes actively misleading.} The nominal $q$ needed to buy the same
$3.0\sigma$-global rate runs $0.38 \to 0.32 \to 1.4\times10^{-2} \to 4.2\times10^{-7}$ across the
rows of Table~\ref{tab:selnoise} (Fig.~\ref{fig:selnoise}, middle). Under coherent defects it
cannot be bought at all: even $q = 10^{-9}$ leaves $2.5\times10^{-3}$ expected false confirmations
at $\epsilon = 10^{-4}$, rising to $0.25$ at $\epsilon = 10^{-2}$, because a coherent artefact
carries a $p$-value small enough to be rejected at essentially any $q$ and then confirms in stage
2. Nothing here contradicts the FDR guarantee --- BH controls the false discovery rate under the
null it is handed, and a contaminated scan does not hand it that null.

\paragraph{The argmax should be dropped rather than deprioritised.} It has no tuning parameter, so
its false-alarm rate is whatever the contamination makes it: $1.35\times10^{-3} \to 0.137 \to
0.687 \to 0.957$ under coherent defects (Fig.~\ref{fig:selnoise}, right). At $\epsilon = 10^{-2}$
the argmax rule produces a false confirmation in $96\,\%$ of scans and there is no way to dial it
back. Under incoherent defects it stays at $1.35\times10^{-3}$, since a glitch dies in stage 2, and
only the power collapses.

\begin{figure}[htbp]
\centering
\includegraphics[width=\textwidth]{bh_outliers.png}
\caption{\textbf{Left:} confirmation probability for a $5\sigma$ signal as the estimator degrades,
all three rules held at the same false-alarm budget; dotted lines are the perfect-estimator
values. \textbf{Middle:} the nominal FDR that buys that budget, against the conventional $q=0.05$.
\textbf{Right:} the argmax's achieved false-alarm rate, which it has no parameter to control.}
\label{fig:selnoise}
\end{figure}

The last row of Table~\ref{tab:selnoise} is the same result as Section~\ref{sec:abnoise} seen from
a different direction: a coherent defect is indistinguishable from a signal, so no selection rule
can help against it. The threshold remains formally calibratable --- it can always be raised, and
its reachable floor is $10^{-9}$ against BH's $2.5\times10^{-3}$ --- but only at essentially zero
power. That is a background-modelling problem, not a selection-rule problem.

Two caveats. The bins here are independent, where a real scan has correlated neighbouring mass
bins; and the defect magnitude is fixed at $2\sigma$, so the size of the gaps scales with it while
the ordering does not. At the heaviest incoherent contamination all three rules are equally dead
and the argmax nominally edges ahead by a few tenths of a point, which is not a meaningful
ranking.

\subsection{What this implies for a data-directed scan}

\begin{enumerate}[leftmargin=1.6em,itemsep=3pt]
  \item \textbf{Do not select by rank alone.} ``Take the most significant bin'' is a dominated
strategy. Selecting on absolute significance recovers the information the ranking throws away, at
identical false-alarm cost, and gains a few percentage points of power at every signal strength.
With an imperfect estimator the argument gets stronger rather than weaker
(Section~\ref{sec:selnoise}): the margin grows to more than ten points, and the argmax loses
control of its false-alarm rate entirely.
  \item \textbf{Fix the fake budget, not the threshold.} The natural, scan-size-independent
        statement of an operating point is ``one expected background candidate per experiment'',
        which determines $t^\star$ from $n$. Quoting a bare $z$ threshold hides the scan size.
  \item \textbf{An FDR rule is a defensible middle ground, but never quote the nominal $q$ as a
false-alarm specification.} BH costs a little power relative to the matched threshold, and it
comes with a guarantee that is easy to state to a reviewer. What that guarantee does not survive
is a miscalibrated tail: the $q$ needed to buy a given false-alarm rate moves by six orders of
magnitude across the contamination range of Section~\ref{sec:selnoise}, and under coherent defects
no $q$ buys it. Whichever rule is adopted, calibrate and quote the \emph{achieved} rate.
  \item \textbf{The confirmation stage is cheap and carries no trials factor.} That is the
        structural fact the next section is built on.
\end{enumerate}
